\item \label{itm:hpmin} {\bf [5 pts.] When does the crust melt?} \\
If temperatures exceed $\sim$850\degC\ in the lower crust, felsic
melts are likely to be generated.
\begin{figure}[hbt]
    \centering
    \includegraphics[]{figures/high_hp.eps}
    \caption{Crustal parameters for question~\ref{itm:hpmin}.}
    \label{fig:highhp}
\end{figure}

What is the minimum heat production these high heat production layers may have
before the mid-crust will melt?
\begin{enumerate}
    \item Develop a relationship for the minimum upper crustal heat
    production before melting occurs.

    \item Why minimum heat production (i.e. consider two regions with the
    same heat flow but different upper crustal heat production)?

    \item Using the following parameters in Table~\ref{tab:highhp}, plot the
    minimum upper crustal heat production as a function of surface heat
    flow, $q_s$, from 20 to 120~mW~m$^{-2}$, typical of continental
    variations.
    \begin{table}[hbt]
        \caption{Parameters for question \ref{itm:hpmin}.}
        \label{tab:highhp}
        \centering
        \begin{tabular}{@{}ll@{}}
            \toprule
            surface temperature & 15$\deg$C \\
            melt temperature & 850$\deg$C \\
            thickness, upper crust & 16~km \\
            melt depth & 30~km \\
            thermal conductivity, upper crust & 2.5 W~m$^{-1}$~K$^{-1}$ \\
            thermal conductivity, lower crust & 2.0 W~m$^{-1}$~K$^{-1}$ \\
            heat production, lower crust & 0.4~$\mu$W~m$^{-3}$ \\
            \bottomrule
        \end{tabular}
    \end{table}

    \item On the same graph, plot the ``average'' upper crustal heat
    production as a function of continental heat flow, $q_r = 0.4 q_s$,
    where $q_r$ is the radiogenic heat flow produced within the upper crust.

    \item In the South Australian heat flow anomaly, heat flow exceeds
    90~mW~m$^{-2}$ and yet there is very little evidence for recent melts.
    Based on your graph, what does this suggest for estimates of upper
    crustal heat production relative to the norm.
\end{enumerate}
\FloatBarrier

